\section{Introduction}
\label{sec:intro}

Chinese ink-wash painting (\emph{Shanshui}, 山水画) is one of the
oldest continuous artistic traditions in the world, governed by
a sophisticated aesthetic theory that is foreign to Western
painting and to the natural-image datasets on which modern
generative models are trained. Two principles are central:
the \emph{five-level ink theory}
(焦、浓、重、淡、清---scorched, thick, heavy, light, clear)
which defines a discrete palette of ink densities, and
\emph{negative space} (留白)---the deliberate use of unpainted
areas as compositional elements, encapsulated in the principle
of \emph{treating white as black} (计白当黑).

Recent advances in image-to-video (I2V) diffusion models such as
Wan2.1-I2V and LTX-Video offer powerful, locally-deployable
backbones for animating static content. However, these models are
trained predominantly on natural imagery and struggle to preserve
the art-specific features of ink-wash paintings: ink-density
distributions are blurred, negative-space regions are filled with
spurious texture, and the temporal motion lacks the gentle
``breathing'' quality of drifting mist or flowing water that
characterizes Shanshui aesthetics.

We propose \textbf{InkFlow}, a dual-conditioned video diffusion
framework that addresses these challenges. The key ideas are:
\begin{enumerate}
    \item \textbf{Dual-stage conditioning}: a line-distribution map
    (the artist's structural draft) is injected during early
    denoising steps to govern composition, while an ink-density
    distribution map (specifying the five ink levels) is injected
    during later steps to govern tonal rendering. This temporal
    separation mirrors the physical ink-wash workflow where
    structure is laid down before tonal detail.
    \item \textbf{Blank-space loss}: a training objective that
    explicitly penalizes texture generation inside user-designated
    negative-space regions, formalizing the
    \emph{treating-white-as-black} principle.
    \item \textbf{Rhythm-of-strokes prior}: a temporal regularizer
    derived from calligraphic timing that encourages smooth,
    breathing-like motion patterns consistent with Shanshui
    aesthetics.
\end{enumerate}

We summarize our contributions as follows:
\begin{enumerate}
    \item A dual-conditioned I2V framework with stage-separated
    injection of line-structure and ink-density conditions, enabling
    art-theoretically grounded control over video generation.
    \item A blank-space loss that preserves negative-space regions
    in generated videos---the first formalization of this
    aesthetic principle as a training objective.
    \item A rhythm-of-strokes temporal prior that produces
    aesthetically appropriate motion patterns for ink-wash
    paintings.
    \item Comprehensive experiments on 30 Shanshui paintings with
    quantitative metrics (ink-density histogram matching,
    blank-space texture ratio, temporal consistency) and a
    perceptual user study.
\end{enumerate}

This work is part of a broader research program on \emph{AI for Art
Process}---using AI to understand and extend the creative
\emph{process} of art-making. Together with work on spatially-aware
stippling and real-time painting coaching, InkFlow contributes to
building AI systems that engage with the \emph{how} of art, not just
the \emph{what}.

The remainder of the paper is organized as follows.
Section~\ref{sec:related} reviews related work in I2V generation,
ink-wash stylization, and controllable diffusion.
Section~\ref{sec:method} presents the InkFlow framework.
Section~\ref{sec:experiments} reports experimental results.
Section~\ref{sec:discussion} discusses limitations and future
directions. Section~\ref{sec:conclusion} concludes.
